\chapter{アカウントをそろえる}
\sheet{2-1〜2-3}

開発で使うサービスのアカウントを用意します。第1巻で必要なのは、GitHub と Claude の2つだけです。
Vercel と Supabase は、使う巻（第3巻）になったら作ります。

\section{GitHub}

\term{GitHub}（ギットハブ）は、コードを保管し、変更の履歴を残し、他の人と共同で作業するためのサービスです。
本番のサイトのコードも、字幕アプリのコードも GitHub にあります。

\begin{step}{GitHub のアカウントを作る}
\begin{enumerate}
  \item ブラウザで \url{https://github.com/signup} を開きます。
  \item メールアドレス、パスワード、ユーザー名を入力します。
  \item 届いたメールの確認コードを入力すると、登録が完了します。
\end{enumerate}

\begin{winwin}[0.7\linewidth]{Sign up to GitHub — Google Chrome}
\centering\textbf{\large Sign up to GitHub}\par\vspace{6pt}
\raggedright
Email*\par
\fbox{\makebox[0.9\linewidth][l]{\color{gray}\scriptsize ryusei.example@gmail.com}}\par\vspace{3pt}
Password*\par
\fbox{\makebox[0.9\linewidth][l]{\color{gray}\scriptsize ••••••••••••}}\par\vspace{3pt}
Username*\par
\fbox{\makebox[0.9\linewidth][l]{\color{gray}\scriptsize shibayama-ryusei}}\par\vspace{6pt}
\centering\colorbox{matcha}{\color{white}\makebox[0.5\linewidth]{Continue}}
\end{winwin}
\figcap{GitHub の登録画面（模式図）}
\end{step}

\begin{caution}[ユーザー名の決め方]
ユーザー名は、URL（\code{github.com/ユーザー名}）にも、コミットの記録にも残ります。
あとから変えると面倒なので、次の点に気をつけて決めてください。
\begin{itemize}
  \item 英小文字・数字・ハイフンだけで作る（例：\code{shibayama-ryusei}、\code{ryusei-s}）
  \item 就職活動などで人に見せても恥ずかしくない名前にする
\end{itemize}
\end{caution}

\begin{step}{2段階認証をオンにする}
GitHub は2段階認証が必須です。登録後しばらくすると設定を求められるので、
スマホの認証アプリ（Google Authenticator や Microsoft Authenticator など）で設定します。
\textbf{リカバリーコードは必ず保存してください}。スマホをなくしたときに、これがないとアカウントに入れなくなります。
\end{step}

\begin{step}{学生向け特典（GitHub Education）を申請する}
\url{https://education.github.com/} から、\textbf{keio.jp のメールアドレス}で学生認証を申請すると、
有料機能の一部が無料になります。必須ではありませんが、通っておくと後で役立ちます。
審査には数日かかることがあります。
\end{step}

\begin{step}{ユーザー名を先輩に伝える}
\sheet{2-2}
練習用のリポジトリにあなたを招待するため、GitHub のユーザー名を LINE で先輩に送ります。
\end{step}

\begin{step}{練習用リポジトリへの招待を承認する}
\sheet{2-2}
先輩が招待すると、GitHub に登録したメールアドレスに
「yutomi141 invited you to yutomi141/rikotaikendo-practice」というメールが届きます。
メールの \btn{View invitation} を押し、開いた画面で \btn{Accept invitation} を押します。
メールが見つからないときは、\url{https://github.com/yutomi141/rikotaikendo-practice/invitations} を直接開いても承認できます。

承認すると、このリポジトリにブランチを作ったり、PR を出したりできるようになります（第2巻から使います）。
\end{step}

\begin{caution}[招待には期限がある]
GitHub の招待は、7日たつと無効になります。期限が切れたら、先輩にもう一度招待してもらってください。
\end{caution}

\section{Claude}

Claude Pro にはすでに入っているので、ここでは2つの点を確認するだけです。

\begin{step}{プランと使用量を確認する}
\begin{enumerate}
  \item \url{https://claude.ai} にログインし、左下のアカウント名から \menu{設定} を開きます。
  \item \menu{使用量}（Usage）の画面で、いまの使用量を確認します。
\end{enumerate}
\end{step}

\begin{tip}[Pro の使用量の仕組み]
\begin{itemize}
  \item 使用量には「5時間ごとの上限」と「週ごとの上限」があります。
  \item \textbf{ブラウザやスマホで Claude とチャットした分と、Claude Code で使った分は、同じ枠から引かれます。}
        開発を進めたい日は、チャットでの雑談を控えめにすると長持ちします。
  \item 上限に達したら、リセットされるまで待ちます。待っている間にできることは第6章で紹介します。
\end{itemize}
\end{tip}

\begin{caution}[追加の課金は押さない]
上限に達すると、追加のクレジットを買うかどうか聞かれることがあります。
\textbf{この教材では課金は一切不要です。}「購入」「有効にする」は押さず、リセットを待ってください。
\end{caution}

\begin{checklist}
\begin{itemize}
  \checkitem GitHub のアカウントを作り、2段階認証を設定した
  \checkitem リカバリーコードを安全な場所に保存した
  \checkitem GitHub のユーザー名を先輩に伝えた
  \checkitem Claude の使用量の画面を開けた
\end{itemize}
\end{checklist}
